
\section{Related Work}

\subsection{Traditional KBQA methods} 
Traditional KBQA methods can be categorized into two primary kinds of approaches: Information Retrieval-based (IR-based) and Semantic Parsing-based (SP-based) ones. IR-based methods~\citep{miller2016kvmem, liu2022graftnet, sun2019pullnet, saxena2020embd, shi2021transfernet, zhang2022subgraph} first retrieve query-relevant information from the KB and subsequently rank candidate answers. Among IR-based methods, some~\citep{bordes2015large, chen2019bidirectional} incorporate a long-term memory component that can be read from and written to, or perform a path-walking process over the KB to derive the answer~\citep{qiu2020stepwise, ren2021lego}. SP-based methods~\citep{sun2020sparqa, chen2019uhop, lanquery, jiangunikgqa, ye2022rng, shu2022tiara} transform the natural language question into a logical query, link the entities and relations, and execute the completed query to obtain the answer. With the advent of LLMs, they have been employed to enhance both approaches, serving as powerful subgraph retrievers~\citep{zhao2024kg, liu2024dual} or semantic parsers~\citep{li2023few, nie2024code}. However, these enhancements do not fundamentally alter the predefined, multi-stage workflow, which brings inflexibility and error propagation.
\subsection{Agentic KBQA methods}
Recently, a new paradigm of agentic reasoning has emerged, which views the LLM as an agent that uses tools to interact with the KB, progressively reason, and finally find the answer. We call them agentic methods in this paper. They can be divided into two main categories: prompt-based and finetuning-based ones. The former methods, such as Interactive-KBQA~\citep{xiong2024interactive}, QueryAgent~\citep{huang2024queryagent}, Readi~\citep{cheng2024call}, and TOG~\citep{sunthink}, leverage carefully designed interaction strategies and prompts to elicit the model's reasoning capabilities. However, these methods struggle to equip the model with sufficient agentic capability to effectively perceive the KB and execute complex reasoning. Their success is also heavily dependent on the intrinsic capabilities of LLMs, most of which are closed-source and have an extremely large number of parameters. However, the limited number of examples is insufficient for the model to fully comprehend complex questions and capture the underlying KB structure, leading to suboptimal performance. Finetuning-based methods aim to distill reasoning abilities into smaller, open-source LLMs. Approaches~\citep{gu2023don,jiang2024kgagent,zhu2025knowagent, fang2024dara} achieve this by constructing high-quality reasoning trajectories and finetuning on them. However, these trajectories, by decomposing and augmenting gold program annotations and the reliance on process supervision via SFT, limit the model's capacity for autonomous exploration, resulting in poor robustness and flexibility. Although some recent works have used Monte Carlo Tree Search (MCTS) methods~\citep{luokbqao1, xiong2025mcts} to expand samples during inference for incremental fine-tuning, this is an indirect way to foster exploration and still operates within the confines of the process-supervised SFT paradigm. Overall, existing finetuning-based methods fall under SFT-based process supervision, which, as noted in Section \ref{sec:intro}, struggles to encourage autonomous exploration, leading to less robust and adaptable reasoning trajectories.